\subsection{分裂半单李代数的根}

在本节中，\((\mathfrak{g},\mathfrak{h})\) 表示一个分裂半单李代数。

对任意 \(\lambda \in \mathfrak{h}^*\) ，用 \(\mathfrak{g}^{\lambda}(\mathfrak{h})\) ，或简称 \(\mathfrak{g}^{\lambda}\) ，表示 \(\mathfrak{g}\) 相对于 \(\lambda\) 的准素子空间（参见第七章 §1 第 3 节）。回顾 \(\mathfrak{g}^{0} = \mathfrak{h}\) （第七章 §2 第 1 节 命题 4），\(\mathfrak{g}\) 是诸 \(\mathfrak{g}^{\lambda}\) 的直和（第七章 §1 第 3 节 命题 8 与 9），\(\mathfrak{g}^{\lambda}\) 是由 \(x \in \mathfrak{g}\) 中使 \([h, x] = \lambda(h)x\) 对一切 \(h \in \mathfrak{h}\) 成立者组成的集合（第七章 §2 第 4 节 定理 2 之推论 1），以及 \(\mathfrak{h}\) 在 \(\mathfrak{g}\) 上的权是线性形式 \(\lambda\) （\(\mathfrak{h}\) 上的），且 \(\mathfrak{g}^{\lambda} \neq 0\) （第七章 §1 第 1 节）。

定义 2. \((\mathfrak{g}, \mathfrak{h})\) 的根是 h 在 g 上的一个非零权。

用 \(R(\mathfrak{g},\mathfrak{h})\) ，或简称 R，表示 \((\mathfrak{g},\mathfrak{h})\) 的根集。我们有

\[
\mathfrak {g} = \mathfrak {h} \oplus \bigoplus_ {\alpha \in R} \mathfrak {g} ^ {\alpha}.
\]

命题 1. 设 \(\alpha, \beta\) 是 \((\mathfrak{g}, \mathfrak{h})\) 的根，并设 \(\langle \cdot, \cdot \rangle\) 是 \(\mathfrak{g}\) 上一个非退化的不变对称双线性型（例如 \(\mathfrak{g}\) 的 Killing 型）。

\begin{enumerate}
\def\labelenumi{(\roman{enumi})}
\item
  若 \(\alpha + \beta \neq 0\) ，则 \(\mathfrak{g}^{\alpha}\) 与 \(\mathfrak{g}^{\beta}\) 正交。\(\langle \cdot, \cdot \rangle\) 在 \(\mathfrak{g}^{\alpha} \times \mathfrak{g}^{-\alpha}\) 上的限制是非退化的。\(\langle \cdot, \cdot \rangle\) 在 \(\mathfrak{h}\) 上的限制是非退化的。
\item
  设 \(x \in \mathfrak{g}^{\alpha}\) ，\(y \in \mathfrak{g}^{-\alpha}\) ，\(h \in \mathfrak{h}\) 。则 \([x, y] \in \mathfrak{h}\) ，并且
\end{enumerate}

\[
\langle h, [ x, y ] \rangle = \alpha (h) \langle x, y \rangle .
\]

断言 (i) 是第七章 §1 第 3 节 命题 10 (iii) 的一个特例。若 \(x \in \mathfrak{g}^{\alpha}\) ，\(y \in \mathfrak{g}^{-\alpha}\) ，\(h \in \mathfrak{h}\) ，我们有 \([x, y] \in \mathfrak{g}^{\alpha - \alpha} = \mathfrak{h}\) ，并且

\[
\langle h, [ x, y ] \rangle = \langle [ h, x ], y \rangle = \langle \alpha (h) x, y \rangle = \alpha (h) \langle x, y \rangle .
\]

定理 1. 设 \(\alpha\) 是 \((\mathfrak{g},\mathfrak{h})\) 的根。

\begin{enumerate}
\def\labelenumi{(\roman{enumi})}
\item
  向量空间 \(\mathfrak{g}^{\alpha}\) 的维数是 1。
\item
  向量子空间 \(\mathfrak{h}_{\alpha} = [\mathfrak{g}^{\alpha},\mathfrak{g}^{-\alpha}]\) 是 \(\mathfrak{h}\) 的向量子空间，其维数为 1。它包含唯一的元素 \(H_{\alpha}\) 使 \(\alpha (H_{\alpha}) = 2\) 。
\item
  向量子空间 \(\mathfrak{s}_{\alpha} = \mathfrak{h}_{\alpha} + \mathfrak{g}^{\alpha} + \mathfrak{g}^{-\alpha}\) 是 \(\mathfrak{g}\) 的李子代数。
\item
  若 \(X_{\alpha}\) 是 \(\mathfrak{g}^{\alpha}\) 的非零元素，则存在唯一的 \(X_{-\alpha} \in \mathfrak{g}^{-\alpha}\) 使 \([X_{\alpha}, X_{-\alpha}] = -H_{\alpha}\) 。设 \(\varphi\) 是从 \(\mathfrak{sl}(2, k)\) 到 \(\mathfrak{g}\) 的线性映射，它把 \(X_{+}\) 映到 \(X_{\alpha}\) ，把 \(X_{-}\) 映到 \(X_{-\alpha}\) ，把 \(H\) 映到 \(H_{\alpha}\) ；则 \(\varphi\) 是从李代数 \(\mathfrak{sl}(2, k)\) 到李代数 \(\mathfrak{s}_{\alpha}\) 的同构。
\end{enumerate}

\begin{enumerate}
\def\labelenumi{\alph{enumi})}
\item
  设 \(h_{\alpha}\) 是 h 的唯一元素，使 \(\alpha(h) = \langle h_{\alpha}, h \rangle\) 对一切 \(h \in h\) 成立。由命题 1，\([x, y] = \langle x, y \rangle h_{\alpha}\) 对一切 \(x \in g^{\alpha}\) ，\(y \in g^{-\alpha}\) 成立；另一方面 \(\langle g^{\alpha}, g^{-\alpha} \rangle \neq 0\) 。因此 \(h_{\alpha} = [g^{\alpha}, g^{-\alpha}] = kh_{\alpha}\) 。
\item
  选取 \(x \in \mathfrak{g}^{\alpha}\) ，\(y \in \mathfrak{g}^{-\alpha}\) 使 \(\langle x, y \rangle = 1\) ，于是 \([x, y] = h_{\alpha}\) 。回顾 \([h_{\alpha}, x] = \alpha(h_{\alpha})x\) ，\([h_{\alpha}, y] = -\alpha(h_{\alpha})y\) 。若 \(\alpha(h_{\alpha}) = 0\) ，则由此推出 \(kx + ky + kh_{\alpha}\) 是幂零子代数 \(t\) （\(\mathfrak{g}\) 中的）；由于 \(h_{\alpha} \in [t, t]\) ，\(\operatorname{ad}_{\mathfrak{g}} h_{\alpha}\) 是幂零的（第一章 §5 第 3 节 定理 1），而这是荒谬的，因为 \(\operatorname{ad}_{\mathfrak{g}} h_{\alpha}\) 非零且半单。所以 \(\alpha(h_{\alpha}) \neq 0\) 。因此存在唯一的 \(H_{\alpha} \in \mathfrak{h}_{\alpha}\) 使 \(\alpha(H_{\alpha}) = 2\) ，这就证明了 (ii)。
\item
  选取非零元素 \(X_{\alpha}\) （\(\mathfrak{g}^{\alpha}\) 中的）。存在 \(X_{-\alpha} \in \mathfrak{g}^{-\alpha}\) 使 \([X_{\alpha}, X_{-\alpha}] = -H_{\alpha}\) （因为 \([X_{\alpha}, \mathfrak{g}^{-\alpha}] = \mathfrak{h}_{\alpha}\) ，由 \(b\) ）。则
\end{enumerate}

\[
\begin{array}{c} [ H _ {\alpha}, X _ {\alpha} ] = \alpha (H _ {\alpha}) X _ {\alpha} = 2 X _ {\alpha}, [ H _ {\alpha}, X _ {- \alpha} ] = - \alpha (H _ {\alpha}) X _ {- \alpha} = - 2 X _ {- \alpha}, \\ \relax [ X _ {\alpha}, X _ {- \alpha} ] = - H _ {\alpha}; \end{array}
\]

于是 \(kX_{\alpha} + kX_{-\alpha} + kH_{\alpha}\) 是 \(\mathfrak{g}\) 的一个子代数，而线性映射 \(\varphi\) （从 \(\mathfrak{sl}(2,k)\) 到 \(kX_{\alpha} + kX_{-\alpha} + kH_{\alpha}\) ，满足 \(\varphi(X_{+}) = X_{\alpha}, \varphi(X_{-}) = X_{-\alpha}, \varphi(H) = H_{\alpha}\) ）是李代数的一个同构。

\begin{enumerate}
\def\labelenumi{\alph{enumi})}
\setcounter{enumi}{3}
\tightlist
\item
  假设 \(\dim g^{\alpha} > 1\) 。设 y 是 \(g^{-\alpha}\) 的一个非零元素。存在非零元素 \(X_{\alpha}\) （\(g_{\alpha}\) 中的）使 \(\langle y, X_{\alpha} \rangle = 0\) 。按 c) 选取 \(X_{-\alpha}\) ，并考虑表示 \(\rho : u \mapsto \operatorname{ad}_{\mathfrak{g}} \varphi(u)\) （从 \(\mathfrak{sl}(2, k)\) 到 g）。我们有
\end{enumerate}

\[
\begin{array}{c} \rho (H) y = [ \varphi (H), y ] = [ H _ {\alpha}, y ] = - 2 y \\ \rho (X _ {+}) y = [ \varphi (X _ {+}), y ] = [ X _ {\alpha}, y ] = \langle X _ {\alpha}, y \rangle h _ {\alpha} = 0. \end{array}
\]

于是，\(y\) 对 \(\rho\) 是本原元，权为 \(-2\) ，这与 §1 第 2 节 命题 2 矛盾。这就证明了 (i)。

\begin{enumerate}
\def\labelenumi{\alph{enumi})}
\setcounter{enumi}{4}
\tightlist
\item
  断言 (iii) 现在是 c) 的推论。另一方面，若 \(X_{\alpha}\) 是 \(g^{\alpha}\) 的非零元素，则在 c) 中构造的元素 \(X_{-\alpha}\) 是 \(g^{-\alpha}\) 中使 \([X_{\alpha}, X_{-\alpha}] = -H_{\alpha}\) 的唯一元素，因为 \(\dim g^{-\alpha} = 1\) 。(iv) 的最后一个断言是 c) 的推论。证毕。
\end{enumerate}

以下将保留记号 \(h_{\alpha}\) ，\(H_{\alpha}\) ，\(s_{\alpha}\) 。（定义 \(h_{\alpha}\) 时，我们取 \(\langle\cdot,\cdot\rangle\) 等于 Killing 型。）若 \(X_{\alpha}\) 是 \(g^{\alpha}\) 的非零元素，则定理 1 的同构 \(\varphi\) 与表示 \(u \mapsto \operatorname{ad}_{\mathfrak{g}} \varphi(u)\) （\(\mathfrak{sl}(2,k)\) 在 g 上的）称为对应于 \(X_{\alpha}\) 的同构与表示。

推论. 设 \(\Phi\) 是 \(\mathfrak{g}\) 的 Killing 型。对一切 \(a, b \in \mathfrak{h}\) ，

\[
\Phi (a, b) = \sum_ {\gamma \in \mathrm{R}} \gamma (a) \gamma (b).
\]

事实上，ad a.ad b 保持每个 \(g^{\gamma}\) 稳定，且它在 \(g^{\gamma}\) 上的限制是比率为 \(\gamma(a)\gamma(b)\) 的位似；若 \(\gamma\neq0\) ，则 \(\dim g^{\gamma}=1\) 。

命题 2. 设 \(\alpha, \beta \in \mathbb{R}\) 。

\begin{enumerate}
\def\labelenumi{(\roman{enumi})}
\item
  \(\beta(H_{\alpha}) \in \mathbf{Z}\) 。
\item
  若 \(\Phi\) 表示 \(\mathfrak{g}\) 的 Killing 型，则 \(\Phi(H_{\alpha}, H_{\beta}) \in \mathbf{Z}\) 。
\end{enumerate}

设 \(X_{\alpha}\) 是 \(\mathfrak{g}^{\alpha}\) 的非零元素，\(\rho\) 是 \(\mathfrak{sl}(2,k)\) 在 \(\mathfrak{g}\) 上对应于 \(X_{\alpha}\) 的表示。\(\rho(\mathrm{H})\) 的特征值是 0 以及诸 \(\beta(H_{\alpha})\)

其中 \(\beta \in \mathbb{R}\) 。因此 (i) 由 §1 第 2 节 命题 2 的推论得出。断言 (ii) 由 (i) 及定理 1 的推论得出。证毕。

设 \(\alpha \in \mathrm{R}\) ，\(X_{\alpha}\) 是 \(\mathfrak{g}^{\alpha}\) 的非零元素，\(X_{-\alpha}\) 是 \(\mathfrak{g}^{-\alpha}\) 中使 \([X_{\alpha}, X_{-\alpha}] = -H_{\alpha}\) 的元素，\(\rho\) 是 \(\mathfrak{sl}(2, k)\) 在 \(\mathfrak{g}\) 上对应于 \(X_{\alpha}\) 的表示。设 \(\pi\) 是 \(\mathbf{SL}(2, k)\) 在 \(\mathfrak{g}\) 上与 \(\rho\) 相容的表示（ \(\S 1\) 第 4 节 定理 2）。由于 \(\operatorname{ad} X_{\alpha}\) 是幂零的（第七章 \(\S 1\) 第 3 节 命题 10 (iv)），\(\pi(e^{X_+}) = e^{\operatorname{ad} X_{\alpha}}\) 是 \(\mathfrak{g}\) 的初等自同构。类似地，\(\pi(e^{X_-}) = e^{\operatorname{ad} X_{-\alpha}}\) 是 \(\mathfrak{g}\) 的初等自同构。因此 \(\pi(\mathbf{SL}(2, k)) \subset \operatorname{Aut}_e(\mathfrak{g})\) 。我们使用记号 \(\theta(t)\) （\(\S 1\) 第 5 节中的）。对 \(t \in k^*\) ，置

\[
\theta_ {\alpha} (t) = \pi (\theta (t)) = e ^ {\mathrm{ad} t X _ {\alpha}} e ^ {\mathrm{ad} t ^ {- 1} X _ {- \alpha}} e ^ {\mathrm{ad} t X _ {\alpha}}.\tag{1}
\]

引理 1. (i) 对一切 \(h \in \mathfrak{h}\) ，\(\theta_{\alpha}(t).h = h - \alpha(h)H_{\alpha}\) 。

\begin{enumerate}
\def\labelenumi{(\roman{enumi})}
\setcounter{enumi}{1}
\item
  对一切 \(\beta \in \mathrm{R}\) ，\(\theta_{\alpha}(t)(\mathfrak{g}^{\beta}) = \mathfrak{g}^{\beta - \beta (H_{\alpha})\alpha}\) 。
\item
  若 \(\alpha, \beta \in \mathrm{R}\) ，则 \(\beta - \beta(H_{\alpha})\alpha \in \mathrm{R}\) 。
\end{enumerate}

设 \(h \in \mathfrak{h}\) 。若 \(\alpha(h) = 0\) ，则 \([X_{\alpha}, h] = [X_{-\alpha}, h] = 0\) ，于是 \(\theta_{\alpha}(t).h = h\) 。另一方面，§1 第 5 节的公式 (5) 表明 \(\theta_{\alpha}(t).H_{\alpha} = -H_{\alpha}\) 。这就证明了断言 (i)。由此得出 \(\theta_{\alpha}(t)^2 | \mathfrak{h} = \mathrm{Id}\) 。若 \(x \in \mathfrak{g}^{\beta}\) 且 \(h \in \mathfrak{h}\) ，

\[
\begin{array}{r l} [ h, \theta_ {\alpha} (t) x ] & = \theta_ {\alpha} (t). [ \theta_ {\alpha} (t) h, x ] - \beta (\theta_ {\alpha} (t) h). \theta_ {\alpha} (t) x \\ & = (\beta (h) - \alpha (h) \beta (H _ {\alpha})). \theta_ {\alpha} (t) x \\ & = (\beta - \beta (H _ {\alpha}) \alpha) (h). \theta_ {\alpha} (t) x \end{array}
\]

于是 \(\theta_{\alpha}(t)x\in \mathfrak{g}^{\beta -\beta (H_{\alpha})\alpha}\) 。这就证明了 (ii)。断言 (iii) 由 (ii) 得出。

定理 2. (i) 集合 \(\mathrm{R} = \mathrm{R}(\mathfrak{g},\mathfrak{h})\) 是 \(\mathfrak{h}^*\) 中的一个既约根系。

\begin{enumerate}
\def\labelenumi{(\roman{enumi})}
\setcounter{enumi}{1}
\tightlist
\item
  设 \(\alpha \in \mathbb{R}\) 。映射 \(s_{\alpha, H_\alpha}: \lambda \mapsto \lambda - \lambda(H_\alpha)\alpha\) （从 \(\mathfrak{h}^*\) 到 \(\mathfrak{h}^*\) ）是唯一的反射 \(s\) （\(\mathfrak{h}^*\) 中的），使 \(s(\alpha) = -\alpha\) 且 \(s(\mathrm{R}) = \mathrm{R}\) 。对一切 \(t \in k^*\) ，\(s\) 是 \(\theta_\alpha(t)|\mathfrak{h}\) 的转置。
\end{enumerate}

首先，R 生成 \(h^{*}\) ，因为若 \(h \in h\) 使 \(\alpha(h) = 0\) 对一切 \(\alpha \in R\) 成立，则 ad h = 0 ，从而由于 g 的中心为零而得 h = 0 。按定义，\(0 \notin R\) 。设 \(\alpha \in R\) 。由于 \(\alpha(H_{\alpha}) = 2\) ，\(s = s_{\alpha,H_{\alpha}}\) 是使 \(s(\alpha) = -\alpha\) 的一个反射。于是由引理 1 (iii) 有 \(s(R) = R\) ，且 \(\beta(H_{\alpha}) \in \mathbf{Z}\) 对一切 \(\beta \in R\) 成立（命题 2 (i)）。这表明 R 是 \(h^{*}\) 中的一个根系。对一切 \(h \in h\) 与 \(\lambda \in h^{*}\) ，

\[
\langle s (\lambda), h \rangle = \langle \lambda - \lambda (H _ {\alpha}) \alpha , h \rangle = \langle \lambda , h - \alpha (h) H _ {\alpha} \rangle = \langle \lambda , \theta_ {\alpha} (t) h \rangle
\]

于是 \(s\) 是 \(\theta_{\alpha}(t)|\mathfrak{h}\) 的转置。最后，我们证明根系 R 是既约的。设 \(\alpha \in \mathrm{R}\) 且 \(y \in \mathfrak{g}^{2\alpha}\) 。由于 \(3\alpha \notin \mathrm{R}\) （第六章 §1 第 3 节 命题 8），\([X_{\alpha}, y] = 0\) ；另一方面，\([X_{-\alpha}, y] \in \mathfrak{g}^{-\alpha + 2\alpha} = \mathfrak{g}^{\alpha} = kX_{\alpha}\) ，故 \([X_{\alpha}, [X_{-\alpha}, y]] = 0\) ；从而

\[
4 y = 2 \alpha (H _ {\alpha}) y = [ H _ {\alpha}, y ] = - [ [ X _ {\alpha}, X _ {- \alpha} ], y ] = 0
\]

于是 \(y = 0\) 且 \(\mathfrak{g}^{2\alpha} = 0\) 。换言之，\(2\alpha\) 不是根。

证毕。

把 \(\mathfrak{h}\) 与 \(\mathfrak{h}^{**}\) 典则地等同起来。采用第六章 §1 第 1 节的记号，由定理 2 (ii) 我们有

\[
H _ {\alpha} = \alpha^ {\vee} \quad \text { for   all } \alpha \in \mathrm{R}.\tag{2}
\]

诸 \(H_{\alpha}\) 于是构成根系 \(\mathbb{R}^{\vee}\) （\(\mathfrak{h}\) 中的），它与 \(\mathbb{R}\) 互逆。

我们称 \(\mathrm{R}(\mathfrak{g},\mathfrak{h})\) 为 \((\mathfrak{g},\mathfrak{h})\) 的根系。反射 \(s_{\alpha,H_{\alpha}}\) 将简记为 \(s_{\alpha}\) 。Weyl 群、权群、Coxeter 数 \ldots{} （\(\mathrm{R}(\mathfrak{g},\mathfrak{h})\) 的）称为 Weyl 群、权群、Coxeter 数 \ldots{} （\((\mathfrak{g},\mathfrak{h})\) 的）。如同第六章 §1 第 1 节，我们把 Weyl 群看作不仅作用于 \(h^{*}\) ，而且由结构转移也作用于 h ，使得 \(s_{\alpha} = \theta_{\alpha}(t)|\mathfrak{h}\) 。由于诸 \(\theta_{\alpha}(t)\) 是 g 的初等自同构，我们有：

推论. \((\mathfrak{g},\mathfrak{h})\) 的 Weyl 群的每个元素在 h 上作用时，都是 g 的某个初等自同构在 h 上的限制。

关于此结果的一个逆命题，见 §5 第 2 节 命题 4。

注 1. 若用 \(\mathfrak{h}_{\mathbf{Q}}\) （相应地 \(\mathfrak{h}_{\mathbf{Q}}^{*}\) ）表示 \(\mathbf{Q}\) -向量空间 \(\mathfrak{h}\) （相应地 \(\mathfrak{h}^{*}\) ）中由诸 \(H_{\alpha}\) （相应地诸 \(\alpha\) ，其中 \(\alpha \in \mathbb{R}\) ）生成的子空间，则 \(\mathfrak{h}\) （相应地 \(\mathfrak{h}^{*}\) ）可以典则地等同于 \(\mathfrak{h}_{\mathbf{Q}} \otimes_{\mathbf{Q}} k\) （相应地 \(\mathfrak{h}_{\mathbf{Q}}^{*} \otimes_{\mathbf{Q}} k\) ），并且 \(\mathfrak{h}_{\mathbf{Q}}^{*}\) 可以等同于 \(\mathfrak{h}_{\mathbf{Q}}\) 的对偶（第六章 §1 第 1 节 命题 1）。我们称 \(\mathfrak{h}_{\mathbf{Q}}\) 与 \(\mathfrak{h}_{\mathbf{Q}}^{*}\) 为 \(\mathbf{Q}\) 在 \(\mathfrak{h}\) 与 \(\mathfrak{h}^{*}\) 上的典则结构（代数 第二章 §8 第 1 节 定义 1）。当提到 \(\mathbf{Q}\) -有理性（对 \(\mathfrak{h}\) 的向量子空间、对 \(\mathfrak{h}\) 上的线性形式等等）时，除非另有说明，我们指的都是这些结构。当提到 R(g, h) 的 Weyl 房或面时，我们将在 \(\mathfrak{h}_{\mathbf{Q}} \otimes_{\mathbf{Q}} \mathbf{R}\) 或 \(\mathfrak{h}_{\mathbf{Q}}^{*} \otimes_{\mathbf{Q}} \mathbf{R}\) 中工作，并把它们记作 \(\mathfrak{h}_{\mathbf{R}}\) 与 \(\mathfrak{h}_{\mathbf{R}}^{*}\) 。

注 2. 根系 \(\mathbb{R}^{\vee}\) （\(\mathfrak{h}\) 中的）定义了一个非退化对称双线性型 \(\beta\) （\(\mathfrak{h}\) 上的）（第六章 §1 第 1 节 命题 3），即形式 \((a,b)\mapsto \sum_{\alpha \in \mathbb{R}}\langle \alpha ,a\rangle \langle \alpha ,b\rangle\) 。由定理 1 的推论，这个形式正是 Killing 型在 \(\mathfrak{h}\) 上的限制。\(\beta |\mathfrak{h}_{\mathbf{Q}}\times \mathfrak{h}_{\mathbf{Q}}\) 到 \(\mathfrak{h}_{\mathbf{Q}}\otimes_{\mathbf{Q}}\mathbf{R}\) 上的扩张是正定的、非退化的（第六章 §1 第 1 节 命题 3）。另一方面，我们可以看到，\(\mathfrak{h}^*\) 上的逆形式（对应于 \(\mathfrak{h}\) 上的一个限制，即 \(\mathfrak{g}\) 的 Killing 型的限制）就是 R 的典则双线性型 \(\Phi_{\mathrm{R}}\) （第六章 §1 第 12 节）。

设 \((\mathfrak{g}_{1},\mathfrak{h}_{1})\) ， \((\mathfrak{g}_{2},\mathfrak{h}_{2})\) 是分裂半单李代数，\(\varphi\) 是从 \(\mathfrak{g}_{1}\) 到 \(\mathfrak{g}_{2}\) 的同构，且 \(\varphi(\mathfrak{h}_{1})=\mathfrak{h}_{2}\) 。由结构转移，映射 \(\varphi|\mathfrak{h}_{1}\) 的转置把 \(R(\mathfrak{g}_{2},\mathfrak{h}_{2})\) 变到 \(R(\mathfrak{g}_{1},\mathfrak{h}_{1})\) 。

命题 3. 设 g 是半单李代数，\(h_{1}\) 与 \(h_{2}\) 是 g 的分裂嘉当子代数。存在从 \(h_{1}^{*}\) 到 \(h_{2}^{*}\) 的同构，它把 \(\mathrm{R}(\mathfrak{g},\mathfrak{h}_{1})\) 变到 \(\mathrm{R}(\mathfrak{g},\mathfrak{h}_{2})\) 。

（更精确的结果，见 §3 第 3 节 命题 10 的推论，以及 §5 第 3 节 命题 5）。

设 \(k'\) 是 \(k\) 的一个代数闭包，\(\mathfrak{g}' = \mathfrak{g} \otimes_k k'\) ，\(\mathfrak{h}_i' = \mathfrak{h}_i \otimes_k k'\) 。那么 \(\mathrm{R}(\mathfrak{g}', \mathfrak{h}_i')\) 是 \(\mathrm{R}(\mathfrak{g}, \mathfrak{h}_i)\) 在映射 \(\lambda \mapsto \lambda \otimes 1\) （从 \(\mathfrak{h}_i^*\) 到 \(\mathfrak{h}_i^* \otimes_k k' = \mathfrak{h}_i'^*\) ）下的像。由第七章 §3 第 2 节 定理 1，存在 \(\mathfrak{g}'\) 的自同构把 \(\mathfrak{h}_1'\) 变到 \(\mathfrak{h}_2'\) ，因而存在同构 \(\varphi\) （从 \(\mathfrak{h}_1'^*\) 到 \(\mathfrak{h}_2'^*\) ），它把 \(\mathrm{R}(\mathfrak{g}', \mathfrak{h}_1')\) 变到 \(\mathrm{R}(\mathfrak{g}', \mathfrak{h}_2')\) 。于是 \(\varphi | \mathfrak{h}_1^*\) 把 \(\mathrm{R}(\mathfrak{g}, \mathfrak{h}_1)\) 变到 \(\mathrm{R}(\mathfrak{g}, \mathfrak{h}_2)\) ，从而把 \(\mathfrak{h}_1^*\) 变到 \(\mathfrak{h}_2^*\) 。证毕。

由命题 3，\((\mathfrak{g},\mathfrak{h})\) 的根系在同构意义下只依赖于 g 而不依赖于 h。同样，Weyl 群、权群 \ldots{} （\((\mathfrak{g},\mathfrak{h})\) 的）按语言滥用简称为 g 的 Weyl 群、权群 \ldots{} （另参见 §5 第 3 节 注 2）。若 g 的邓金图是 \(A_{l}\) 型，或 \(B_{l}\) 型，\ldots{} （参见第六章 §4 第 2 节 定理 3），我们就说 g 是 \(A_{l}\) 型，或 \(B_{l}\) 型，\ldots 型的。

回顾，若 \(\alpha\) 与 \(\beta\) 是线性无关的根，则 \(j \in \mathbf{Z}\) 中使 \(\beta + j\alpha \in \mathrm{R}\) 者的集合是区间 \([-q, p]\) （\(\mathbf{Z}\) 中的，包含 0），且 \(p - q = -\langle \beta, \alpha^{\vee} \rangle = -\beta(H_{\alpha})\) （第六章 §1 第 3 节 命题 9）。

命题 4. 设 \(\alpha\) 与 \(\beta\) 是线性无关的根。设 \(p\) （相应地 \(q\) ）是最大的整数 \(j\) ，它使 \(\beta + j\alpha\) （相应地 \(\beta - j\alpha\) ）为根。

\begin{enumerate}
\def\labelenumi{(\roman{enumi})}
\item
  向量子空间 \(\sum_{-q\leq j\leq p}\mathfrak{g}^{\beta +j\alpha}\) （\(\mathfrak{g}\) 中的）是单 \(\mathfrak{s}_{\alpha}\) -模，维数为 \(p + q + 1\) 。
\item
  若 \(\alpha + \beta\) 是根，则 \([\mathfrak{g}^{\alpha}, \mathfrak{g}^{\beta}] = \mathfrak{g}^{\alpha + \beta}\) 。
\end{enumerate}

设 \(X_{\alpha}\) （相应地 x）是 \(\mathfrak{g}^{\alpha}\) （相应地 \(\mathfrak{g}^{\beta+p\alpha}\) ）的非零元素。则

\[
\begin{array}{c} {[ X _ {\alpha}, x ] \in \mathfrak {g} ^ {\beta + (p + 1) \alpha} = 0} \\ {[ H _ {\alpha}, x ] = (\beta (H _ {\alpha}) + p \alpha (H _ {\alpha})) x = (- p + q + 2 p) x = (p + q) x.} \end{array}
\]

于是，\(x\) 是权为 \(p + q\) 的本原元（对 \(\mathfrak{sl}(2,k)\) 在 \(\mathfrak{g}\) 上对应于 \(X_{\alpha}\) 的表示而言）；而 \(\mathfrak{sl}(2,k)\) -模 \(\sum_{-q\leq j\leq p}\mathfrak{g}^{\beta +j\alpha}\) 的维数为 \(p + q + 1\) ；因此它是单模（§1 第 2 节 命题 2）。若 \(\alpha +\beta \in \mathrm{R}\) ，则 \(p\geq 1\) ，于是 \(\mathfrak{g}^{\beta}\) 的元素不是本原的，从而 \([X_{\alpha},\mathfrak{g}^{\beta}]\neq 0\) 。由于 \([\mathfrak{g}^{\alpha},\mathfrak{g}^{\beta}]\subset \mathfrak{g}^{\alpha +\beta}\) ，我们最终得到 \([\mathfrak{g}^{\alpha},\mathfrak{g}^{\beta}] = \mathfrak{g}^{\alpha +\beta}\) 。

注 3. 回顾，由第六章 §1 第 3 节 命题 9 的推论，整数 \(p + q + 1\) 只能取值 1, 2, 3, 4。

注 4. 设 \((\mathfrak{g},\mathfrak{h})\) 是分裂约化李代数，\(\mathfrak{c}\) 是 \(\mathfrak{g},\mathfrak{g}' = \mathcal{D}\mathfrak{g}\) 中 g 的中心，\(\mathfrak{h}' = \mathfrak{h} \cap \mathfrak{g}'\) 。则 \(\mathfrak{h} = \mathfrak{c} \times \mathfrak{h}'\) ，并且我们把 \(\mathfrak{h}'^*\) 等同于 \(\mathfrak{h}^*\) 的一个向量子空间。对任意 \(\lambda \in \mathfrak{h}^*\) ，\(\lambda \neq 0\) ，准素子空间 \(\mathfrak{g}^\lambda\) （相对于 \(\lambda\) 的）等于 \(\mathfrak{g}'^{\lambda | \mathfrak{h}'}\) 。\(\mathfrak{h}\) 在 \(\mathfrak{g}\) 上的非零权称为 \((\mathfrak{g},\mathfrak{h})\) 的根；每个根在 \(\mathfrak{c}\) 上为零。用 \(\mathrm{R}(\mathfrak{g},\mathfrak{h})\) 表示 \((\mathfrak{g},\mathfrak{h})\) 的根集；它可以典则地等同于 \(\mathrm{R}(\mathfrak{g}',\mathfrak{h}')\) 。设 \(\alpha \in \mathrm{R}(\mathfrak{g},\mathfrak{h})\) 。我们如同半单情形那样定义 \(h_\alpha\) ，\(H_\alpha\) ，\(\mathfrak{s}_\alpha\) ，同构 \(\mathfrak{sl}(2,k) \to \mathfrak{s}_\alpha\) ，以及 \(\mathfrak{sl}(2,k)\) 在 \(\mathfrak{g}\) 上对应于 \(\alpha\) 的表示。

